\chapter{Alkaline Phosphatase (ALP)}

\section*{What Is This Test?}
Alkaline phosphatase (ALP) is a group of isoenzymes that catalyze the hydrolysis of phosphate esters at an alkaline pH (optimum pH 9--10). ALP is present in virtually all body tissues, with the highest concentrations found in the liver (canalicular membrane of hepatocytes), bone (osteoblasts), intestine (brush border), placenta (syncytiotrophoblast), and kidney (proximal tubule). In healthy adults, serum ALP is derived primarily from liver and bone in approximately equal proportions. The liver isoenzyme is located on the canalicular membrane of hepatocytes and is released into the blood when there is cholestasis (bile flow obstruction) or hepatic infiltration. The bone isoenzyme reflects osteoblast activity and is elevated in states of increased bone formation (growing children, healing fractures, Paget disease, osteoblastic bone metastases). Intestinal ALP rises after meals (especially fatty meals) in individuals with blood type B or O (secretor status) \cite{schiele1988total}. Placental ALP rises during the third trimester of pregnancy. The half-life of ALP is approximately 7 days. ALP is a core component of the comprehensive metabolic panel and liver function tests. When ALP is elevated, the source must be determined: if GGT or 5'-nucleotidase is also elevated, the ALP elevation is of hepatobiliary origin; if GGT is normal, the ALP elevation is of bone or other non-hepatic origin \cite{kwo2017acg}. ALP isoenzyme fractionation by electrophoresis can definitively identify the tissue source but is rarely needed.

\section*{Alternative Names}
ALP; Alk Phos; Serum Alkaline Phosphatase; Alkaline Phosphatase Isoenzymes; Bone-Specific Alkaline Phosphatase (BSAP); Liver Alkaline Phosphatase

\section*{Principle}
ALP is measured by a kinetic spectrophotometric method using p-nitrophenylphosphate (pNPP) as the substrate. In alkaline buffer (pH 10.2--10.4), ALP cleaves phosphate from pNPP, releasing p-nitrophenol, which is yellow and absorbs light at 405--410 nm: p-Nitrophenylphosphate + H\textsubscript{2}O $\rightarrow$ p-Nitrophenol + Phosphate. The rate of increase in absorbance at 405 nm is directly proportional to ALP activity. The IFCC standardized method uses 2-amino-2-methyl-1-propanol (AMP) buffer with magnesium and zinc as activators (ALP requires zinc and magnesium as cofactors). Results are reported in IU/L at 37 degrees C. ALP isoenzyme fractionation uses heat stability testing (bone ALP is heat-labile, denatured at 56 degrees C for 10 minutes; liver ALP is more heat-stable; placental ALP is most heat-stable, surviving 65 degrees C for 30 minutes), chemical inhibition (L-phenylalanine inhibits intestinal and placental ALP; L-homoarginine inhibits liver and bone ALP), and electrophoresis. The test is performed on automated chemistry analyzers using spectrophotometric methods, not by hematology analyzers.

\section*{History}
Alkaline phosphatase was first discovered in 1907 by Suzuki, Yoshimura, and Takaishi. The role of ALP in bone metabolism was recognized by Robert Robison in 1923, who proposed that ALP was essential for bone mineralization \cite{robison1923possible}. The clinical use of ALP for detecting liver and bone disease developed in the 1930s--1950s. The Bessey-Lowry method using p-nitrophenylphosphate was developed in 1946 and remains the basis of modern ALP assays \cite{bessey1946method}. Heat fractionation to differentiate liver from bone ALP was developed in the 1960s. The IFCC standardized method was published in 2011 \cite{ifcc2011ifcc}.

\section*{Why This Test Is Ordered}
\begin{itemize}
  \item Screening for cholestasis and biliary obstruction as part of the CMP and LFTs
  \item Evaluation of suspected biliary obstruction from gallstones (choledocholithiasis), strictures, primary sclerosing cholangitis (PSC), or pancreatic head malignancy
  \item Diagnosis and monitoring of primary biliary cholangitis (PBC): ALP elevation (often >2--3$\times$ ULN) is a hallmark, confirmed by positive anti-mitochondrial antibody (AMA) \cite{lindor2019primary}
  \item Evaluation of infiltrative liver disease: sarcoidosis, tuberculosis, fungal infections, amyloidosis, lymphoma, metastatic malignancy (ALP may be disproportionately elevated relative to AST/ALT)
  \item Assessment of bone metabolism: Paget disease of bone (markedly elevated ALP with normal GGT), osteomalacia, rickets, bone metastases (especially osteoblastic from prostate cancer), primary bone tumors (osteosarcoma), hyperparathyroidism, healing fractures
  \item Physiological elevation in normal growth: children and adolescents have higher ALP due to active bone growth (may be 1.5--3$\times$ adult upper limit)
  \item Monitoring of antiresorptive therapy (bisphosphonates, denosumab) in Paget disease
\end{itemize}

\section*{Primary vs Secondary Test}
ALP is a primary screening test as part of the CMP. When elevated, GGT or 5'-nucleotidase confirms hepatobiliary origin. It is a first-line marker for cholestasis and bone disease.

\section*{How This Test Is Performed}
A healthcare professional draws blood from a vein into a serum separator tube (gold-top) or lithium heparin tube (green-top). EDTA or citrate tubes must NOT be used because they chelate zinc and magnesium, which are essential cofactors for ALP activity, producing falsely low results. Hemolyzed samples are acceptable for ALP (unlike AST and LDH). The sample is centrifuged and ALP measured on an automated chemistry analyzer. Results are available within 30--60 minutes.

\section*{What Preparation Is Needed}
Fasting for 8--12 hours is recommended. A fatty meal increases intestinal ALP release in individuals with blood type B or O (secretors), which can increase total ALP by 10--20\% for 2--12 hours after the meal. Fasting eliminates this variability. Inform the provider of: (1) Medications that increase ALP: anticonvulsants (phenytoin, phenobarbital---increase ALP synthesis through hepatic enzyme induction), estrogens/oral contraceptives, androgens, erythromycin, chlorpromazine, allopurinol. (2) Medications that decrease bone ALP: bisphosphonates, denosumab, calcitonin. (3) Albumin infusion: some albumin preparations contain ALP. (4) Recent fractures: ALP rises during healing (peak at 2--4 weeks post-fracture). (5) Pregnancy: placental ALP increases in the third trimester, sometimes markedly. Intestinal ALP also increases.

\section*{Contraindications and Precautions}
No absolute contraindications. Important considerations: (1) Zinc and magnesium chelators (EDTA, citrate, oxalate) in collection tubes falsely low ALP because these metals are required cofactors. Only serum or heparin tubes should be used. (2) ALP is elevated in normal growth---children and adolescents have higher ALP than adults; reference ranges for pediatric patients are age-specific. (3) Benign transient hyperphosphatasemia of infancy/childhood: striking elevation of ALP (3--20$\times$ ULN) without liver or bone disease, resolving spontaneously within weeks to months. Usually due to increased sialylation of ALP reducing hepatic clearance. ALP isoenzymes show bone and liver fractions. No treatment needed beyond reassurance. (4) ALP rises during third trimester of pregnancy from placental isoenzyme (1.5--3$\times$ ULN in 30--40\% of pregnancies). (5) Blood group B or O individuals (secretors) have higher intestinal ALP after meals; fasting sample avoids this. (6) Macro-ALP: rare complex of ALP with immunoglobulins causing persistently elevated ALP; differentiated by PEG precipitation \cite{wallach2022interpretation}.

\section*{Risks}
Minimal risk from venipuncture. Mild pain or bruising at site resolves in 1--2 days.

\section*{What Happens After the Test}
Results available within 30--60 minutes. When ALP is elevated, determine hepatobiliary vs. bone origin: (1) Check GGT: if GGT elevated, ALP is heptobiliary (cholestasis). (2) If GGT normal, ALP is likely bone (Paget disease, fractures, bone metastases, osteomalacia, growing child) or placental. (3) If both ALP and GGT are elevated but AST/ALT are disproportionately high, this may indicate a mixed hepatocellular-cholestatic pattern. Biliary obstruction typically causes ALP >3--5$\times$ ULN with GGT >3$\times$ ULN and minimal AST/ALT elevation (or only mild elevation). Infiltrative liver disease (metastases, granulomas) can produce isolated or disproportionate ALP elevation.

\section*{Understanding Results}

\subsection*{Below Normal}
\begin{itemize}
  \item \textbf{Hypophosphatasia:} Rare inherited disorder (autosomal recessive or dominant, ALPL gene mutation) of defective bone mineralization characterized by persistently low ALP, elevated urinary phosphoethanolamine, and bone fragility (rickets in children, osteomalacia in adults, premature tooth loss). Perinatal form is often fatal; milder forms present in adulthood with fractures \cite{whyte2016hypophosphatasia}.
  \item \textbf{Zinc deficiency:} Zinc is a required cofactor for ALP. Severe zinc deficiency (acrodermatitis enteropathica, TPN without zinc, alcoholism) causes low ALP. ALP rises with zinc replacement.
  \item \textbf{Magnesium deficiency:} Magnesium is also a cofactor; severe deficiency can lower ALP.
  \item \textbf{Hypothyroidism:} Reduced osteoblast activity lowers bone ALP.
  \item \textbf{Pernicious anemia / Vitamin B12 deficiency:} Can be associated with low ALP (mechanism unclear).
  \item \textbf{Malnutrition, scurvy (vitamin C deficiency):} Reduced bone formation.
  \item \textbf{Wilson disease:} Associated with low ALP, especially when presenting with fulminant hepatic failure and hemolysis.
  \item \textbf{EDTA/citrate tube contamination:} Falsely low due to zinc chelation.
\end{itemize}

\subsection*{Above Normal}
\begin{itemize}
  \item \textbf{Hepatobiliary causes (elevated GGT):} Biliary obstruction: choledocholithiasis, malignant obstruction (pancreatic head, cholangiocarcinoma, ampullary carcinoma), benign strictures, primary sclerosing cholangitis (ALP frequently >3$\times$ ULN), biliary atresia. Primary biliary cholangitis (ALP >2--3$\times$ ULN, AMA positive, predominantly middle-aged women, associated with fatigue and pruritus). Drug-induced cholestasis (estrogens, anabolic steroids, erythromycin, chlorpromazine, amoxicillin-clavulanate, azathioprine, cyclosporine). Infiltrative disease: metastases, lymphoma, sarcoidosis, amyloidosis, TB. Hepatic congestion (right-sided heart failure, Budd-Chiari syndrome). Sepsis.
  \item \textbf{Bone causes (normal GGT):} Paget disease of bone (markedly elevated ALP---often 3--20$\times$ ULN with normal GGT; focal disorder of accelerated bone remodeling causing bone pain, deformity, fractures, deafness, high-output cardiac failure) \cite{ralston2008pathogenesis}. Bone metastases (osteoblastic---prostate cancer most commonly; also breast, lung). Osteomalacia/rickets (vitamin D deficiency). Primary hyperparathyroidism (mild ALP elevation from increased bone turnover, approximately 10--30\% of patients). Hyperthyroidism (increased bone turnover). Fracture healing (ALP rises 2--4 weeks post-fracture, peaks at 4--8 weeks). Osteosarcoma. Renal osteodystrophy (CKD-MBD---increased bone turnover from secondary hyperparathyroidism; ALP often >2$\times$ ULN).
  \item \textbf{Physiological:} Children and adolescents (bone growth---ALP 1.5--3$\times$ adult ULN). Pregnancy (third trimester placental ALP, 1.5--3$\times$ ULN). Healing fractures. Postmenopausal women (mildly elevated bone ALP).
  \item \textbf{Others:} Benign transient hyperphosphatasemia of infancy/childhood (marked ALP elevation, 3--20$\times$ ULN, self-limited). Certain cancers producing Regan isoenzyme (placental-like ALP, e.g., ovarian, lung, Hodgkin lymphoma). "
\end{itemize}

\section*{Normal Results}
\begin{itemize}
  \item \textbf{Adults:} 44--147 IU/L (varies by method and laboratory; 30--120 IU/L common)
  \item \textbf{Children (1--12 years):} 100--350 IU/L (higher due to bone growth)
  \item \textbf{Adolescents (12--18 years):} 50--260 IU/L (pubertal growth spurt)
  \item \textbf{Infants (<1 year):} 60--300 IU/L
  \item \textbf{Pregnancy (third trimester):} Up to 2--3$\times$ ULN (placental ALP)
  \item \textbf{Elderly:} May have mildly higher ALP than young adults (increased bone turnover)
\end{itemize}

\section*{Follow-up Tests}
\begin{itemize}
  \item \textbf{GGT or 5'-nucleotidase:} If elevated, confirms hepatobiliary origin of ALP. If normal, suggests bone or other non-hepatic source.
  \item \textbf{ALP isoenzyme fractionation (electrophoresis):} Definitive identification of tissue source when origin is unclear.
  \item \textbf{Abdominal ultrasound:} First-line for biliary obstruction, gallstones, dilated ducts.
  \item \textbf{MRCP or ERCP:} Detailed biliary imaging for obstruction, strictures, stones.
  \item \textbf{Anti-mitochondrial antibody (AMA):} Diagnostic for primary biliary cholangitis (positive in >95\%).
  \item \textbf{Bone-specific alkaline phosphatase (BSAP):} More specific marker of bone formation than total ALP.
  \item \textbf{Bone scan (nuclear medicine):} For Paget disease, bone metastases, fractures.
  \item \textbf{Serum calcium, phosphate, PTH, 25(OH)D:} For metabolic bone disease workup.
  \item \textbf{Prostate-specific antigen (PSA):} If prostate cancer metastases suspected.
\end{itemize}

\section*{References}
\bibliographystyle{unsrtnat}
\bibliography{references}

\clearpage
\section*{Regulatory Status and Authorization Date}
Regulatory status applies to the specific assay, instrument, manufacturer, and intended use rather than to the test name alone. No single approval date can be assigned to this test category. For a commercial assay, verify the current product in the relevant regulator's database: the U.S. Food and Drug Administration (FDA) clearance, approval, or authorization; India's Central Drugs Standard Control Organisation (CDSCO) licence or permission; the European Union In Vitro Diagnostic Regulation (EU IVDR) CE certificate issued through the applicable conformity-assessment route; or the United Kingdom Medicines and Healthcare products Regulatory Agency (MHRA) registration and UKCA/CE status. Laboratory-developed tests may not have an individual FDA, CDSCO, EU IVDR, or MHRA approval date.
\clearpage
